\section{Methodology}
\label{sec:methodology}

\subsection{System Architecture and Operational Flow}
\label{sec:architecture}

The proposed framework is implemented as a set of loosely coupled services designed to generate and validate 4D trajectories for structured U-space corridor operations. The main logical components are:
(i) a \emph{Topology Registry} (Webserver) providing vertiports and the corridor network,
(ii) the \texttt{acc-trajcalc} \emph{Path Planner} responsible for route computation and CFMS-based 4D trajectory generation,
(iii) the \texttt{acc-trajdeconf} \emph{Deconfliction Service} responsible for conflict detection and resolution,
and (iv) a \emph{U-space Service Provider (USSP)} responsible for flight-plan validation, authorization, and supervision. NDMap is not part of \texttt{acc-trajcalc}; NDMap-based conflict assessment is encapsulated in the separate \texttt{acc-trajdeconf} service.

\begin{figure}[!htbp]
  \centering
  \begin{tikzpicture}[
    box/.style={
      draw,
      rounded corners=10pt,
      align=center,
      minimum width=38mm,
      minimum height=18mm
    },
    arrow/.style={-Latex, thick},
    node distance=18mm and 38mm
  ]

  % Nodes (2x2 grid)
  \node[box] (web) {Webserver\\(Topology Registry)};
  \node[box, right=of web] (ussp) {U-space Service\\Provider};

  \node[box] (planner) [below=of web] {\texttt{acc-trajcalc}\\(A* + CFMS)};
  \node[box] (deconf)  [below=of ussp] {\texttt{acc-trajdeconf}\\(NDMap + resolver)};

  % Arrows
  % Webserver -> USSP (top horizontal)
  \draw[arrow] (web.east) -- (ussp.west);

  % Webserver -> Planner (left vertical) with label
  \draw[arrow] (web.south) -- node[left, align=center] {vertiports,\\corridor network} (planner.north);

  % Planner -> Deconfliction service (bottom horizontal) with label
  \draw[arrow] (planner.east) -- node[above, align=center] {candidate 4D\\trajectory} (deconf.west);

  % Deconfliction service -> USSP (right vertical) with label
  \draw[arrow] (deconf.north) -- node[right, align=center] {conflict-free\\trajectory or rejection} (ussp.south);

  \end{tikzpicture}
  \caption{High-level service architecture separating trajectory calculation from NDMap-based deconfliction.}
  \label{fig:architecture}
\end{figure}

Fig.~\ref{fig:architecture} illustrates the high-level service interactions. The detailed internal planning workflow, including iterative conflict detection
and lateral re-routing in the deconfliction service, is depicted in Fig.~\ref{fig:planning_loop}.

\subsection{Global Route Planning on Corridor Graphs}
\label{sec:global_planning}

Global routing between vertiports is formulated as a shortest-path problem on a weighted graph \(G=(V,E)\),
where nodes \(V\) represent vertiports and edges \(E\) represent available corridor connections.
Edge weights correspond to geodesic distance between the connected endpoints (Haversine distance).
A* search is used to compute the shortest route between the selected origin and destination.
The heuristic is the straight-line geodesic distance to the destination, which is admissible and yields efficient convergence.
This heuristic is admissible because any feasible route constrained to graph edges cannot be shorter than the direct geodesic distance from the current node to the destination.
The output of this stage is an ordered list of corridor nodes:
\[
P = [v_0, v_1, \dots, v_n],
\]
which serves as the lateral (2D) reference route for subsequent trajectory generation.

% Optional (only if you want it in the paper):
% Copy your thesis A* illustration into figures/astar_graph.png and uncomment:
% \begin{figure}[!htbp]
%   \centering
%   \includegraphics[width=0.70\linewidth]{figures/astar_graph.png}
%   \caption{Illustration of A* search on the corridor graph.}
%   \label{fig:astar_graph}
% \end{figure}

\subsection{4D Trajectory Generation via CFMS}
\label{sec:trajectory_generation}

The route \(P\) is transformed into a phase-consistent four-dimensional trajectory (latitude, longitude, altitude, time) using a waypoint synthesis step and the CFMS flight-management and 4D trajectory generation system~\cite{morscheck2021cfms}.

\begin{algorithm}[htbp]
\caption{Planner pipeline for 4D trajectory generation}
\label{alg:planner_pipeline}
\begin{algorithmic}[1]
\REQUIRE Origin vertiport \(o\), destination vertiport \(d\), corridor graph \(G=(V,E)\), vertiport data, cruise altitude \(H_c\)
\ENSURE Dense trajectory \(\mathcal{T}_{\mathrm{dense}}\) and subsampled output trajectory \(\mathcal{T}_{\mathrm{out}}\)
\STATE Compute corridor route \(P\) from \(o\) to \(d\) using A* on \(G\)
\STATE Generate phase-structured waypoint sequence \(W\) from \(P\), vertiport data, and \(H_c\)
\STATE Convert \(W\) into CFMS-compatible input
\STATE Generate dense 4D trajectory \(\mathcal{T}_{\mathrm{dense}}\) using CFMS
\STATE Create subsampled output trajectory \(\mathcal{T}_{\mathrm{out}}\) with at most 200 trajectory points
\RETURN \(\mathcal{T}_{\mathrm{dense}}, \mathcal{T}_{\mathrm{out}}\)
\end{algorithmic}
\end{algorithm}

Fig.~\ref{fig:waypoint-synthesis-schematic} summarizes the waypoint-synthesis
logic used between A* routing and CFMS execution. A* provides the lateral
corridor-node sequence, while the waypoint synthesizer inserts terminal
waypoints around the departure and arrival vertiports. Intermediate A* corridor
nodes are retained as cruise-route references and marked as navigation nodes in
the waypoint list. The resulting phase-structured waypoint sequence is then
converted into CFMS input for dense 4D trajectory generation.

\begin{figure}[!htbp]
\centering
\resizebox{\linewidth}{!}{%
\begin{tikzpicture}[
    >=Latex,
    font=\small,
    route/.style={very thick, blue!70!black},
    waypoint/.style={circle, draw=black, fill=white, inner sep=1.6pt},
    aux/.style={gray!55, dashed, thin},
    leader/.style={gray!65, thin},
    labelbox/.style={font=\scriptsize, align=center, fill=white, inner sep=1.8pt},
    smalllabel/.style={font=\scriptsize, align=center, fill=white, inner sep=1.5pt},
    volume/.style={fill=blue!10, draw=blue!35, dashed, rounded corners=2pt},
    axis/.style={gray!65, thin, -Latex}
]

% ------------------------------------------------------------------
% Coordinates
% ------------------------------------------------------------------
\coordinate (origin) at (0,0);
\coordinate (depHone) at (0,1.05);
\coordinate (depHtwo) at (1.55,2.05);
\coordinate (ols) at (2.65,3.15);
\coordinate (cruiseA) at (3.55,4.20);
\coordinate (cruiseB) at (5.80,4.20);
\coordinate (cruiseC) at (8.05,4.20);
\coordinate (tod) at (9.35,4.20);
\coordinate (arrEntry) at (10.65,3.05);
\coordinate (arrHtwo) at (11.75,2.05);
\coordinate (arrHone) at (11.75,1.05);
\coordinate (ldp) at (11.75,0.55);
\coordinate (dest) at (11.75,0);

% ------------------------------------------------------------------
% Background reference lines and axes
% ------------------------------------------------------------------
\draw[axis] (-0.85,-0.35) -- (13.10,-0.35);
\draw[axis] (-0.85,-0.35) -- (-0.85,5.05);

\node[rotate=90, gray!70, font=\scriptsize] at (-1.25,2.45) {altitude / MSL};
\node[gray!70, font=\scriptsize] at (6.15,-1.15)
{progress along selected A* corridor route};

\draw[aux] (-0.55,4.20) -- (12.75,4.20);
\node[labelbox, anchor=west] at (10.15,4.45) {$H_c = 200\,\mathrm{m}$ MSL};

\draw[aux] (-0.55,2.05) -- (0.60,2.05);
\draw[aux] (11.45,2.05) -- (12.95,2.05);
\node[smalllabel, font=\small, text=black, anchor=west] at (-0.48,2.22) {$h_{2}$};

\draw[aux] (-0.55,1.05) -- (0.45,1.05);
\draw[aux] (11.45,1.05) -- (12.95,1.05);
\node[smalllabel, font=\small, text=black, anchor=west] at (-0.48,1.22) {$h_{1}$};

\draw[aux] (-0.55,0.55) -- (0.45,0.55);
\draw[aux] (11.45,0.55) -- (12.95,0.55);

% ------------------------------------------------------------------
% Terminal OFV / verticone volumes
% ------------------------------------------------------------------
\begin{scope}[on background layer]
    \path[volume]
        (-0.18,0.00) --
        (0.62,0.00) --
        (1.05,1.45) --
        (1.80,2.55) --
        (1.35,2.80) --
        (-0.18,2.80) -- cycle;

    \path[volume]
        (11.15,0.00) --
        (12.30,0.00) --
        (12.30,2.80) --
        (10.75,2.80) --
        (11.25,2.10) --
        (11.55,1.10) -- cycle;
\end{scope}

% ------------------------------------------------------------------
% Main synthesized waypoint route
% ------------------------------------------------------------------
\draw[route]
    (origin) --
    (depHone) --
    (depHtwo) --
    (ols) --
    (cruiseA) --
    (cruiseB) --
    (cruiseC) --
    (tod) --
    (arrEntry) --
    (arrHtwo) --
    (arrHone) --
    (ldp) --
    (dest);

% Waypoint markers
\foreach \p in {origin,depHone,depHtwo,ols,cruiseA,cruiseB,cruiseC,tod,arrEntry,arrHtwo,arrHone,ldp,dest}
{
    \node[waypoint] at (\p) {};
}

% ------------------------------------------------------------------
% Labels: departure side
% ------------------------------------------------------------------
\node[labelbox, anchor=north] at (0,-0.48) {origin pad\\origin};

\node[labelbox, anchor=east] at (-0.45,0.72) {vertical\\climb};
\draw[leader] (-0.10,0.72) -- ($(origin)!0.55!(depHone)$);

\node[labelbox] at (0.78,3.30) {departure OFV\\verticone};

\node[labelbox] at (2.35,1.48) {transition};
\draw[leader] (2.05,1.58) -- ($(depHone)!0.62!(depHtwo)$);

\node[labelbox] (olslabel) at (2.35,4.95) {optional\\OLS point};
\draw[leader] (olslabel.south) to[out=-90,in=120] (ols);

% ------------------------------------------------------------------
% Labels: cruise / corridor middle
% ------------------------------------------------------------------
\node[labelbox] at (5.95,5.02) {intermediate corridor nodes\\cruise references};
\draw[leader] (5.80,4.78) -- (cruiseB);
\draw[leader] (8.05,4.78) -- (cruiseC);

\node[labelbox] at (9.35,5.02) {top of\\descent};
\draw[leader] (9.35,4.78) -- (tod);

\node[
    draw=gray!55,
    rounded corners=2pt,
    fill=white,
    align=center,
    inner sep=4pt
] at (6.20,2.45)
{waypoint synthesis adds\\vertical and terminal structure\\[2pt]
 CFMS-compatible waypoint sequence\\[-1pt]
 \(\displaystyle W=[w_1,\ldots,w_m]\)};

\node[smalllabel] at (5.90,-0.58)
{\(v_1,\ldots,v_n\): A* corridor-node sequence \(P\)};

% Light gray guide dots for corridor-node sequence
\foreach \x in {1.7,3.55,5.80,8.05,9.35,10.65}
{
    \node[circle, fill=gray!45, inner sep=1.1pt] at (\x,-0.35) {};
}

% ------------------------------------------------------------------
% Labels: arrival side
% ------------------------------------------------------------------
\node[labelbox] at (11.80,3.52) {arrival OFV\\verticone};

\node[labelbox, anchor=west] at (12.55,2.38) {OFV entry};
\draw[leader] (12.35,2.30) -- (arrHtwo);

\node[labelbox, anchor=west] at (12.55,1.40) {over-pad\\waypoint};
\draw[leader] (12.35,1.18) -- (arrHone);

\node[labelbox, anchor=west] at (12.55,0.58) {LDP};
\draw[leader] (12.30,0.58) -- (ldp);

\node[labelbox, anchor=north] at (11.75,-0.48) {destination pad\\destination};

\end{tikzpicture}%
}
\caption{Schematic waypoint-synthesis procedure used to convert an A*-selected corridor route into a CFMS-compatible phase-structured waypoint sequence. The departure and arrival procedures introduce terminal waypoints at \(h_1\), \(h_2\), and the LDP, while intermediate corridor nodes define cruise-route references at the fixed cruise layer.}
\label{fig:waypoint-synthesis-schematic}
\end{figure}

\paragraph{Inputs.}
The trajectory generation stage consumes:
(i) the node sequence \(P\),
(ii) vertiport coordinates and elevations,
(iii) a fixed cruise altitude layer, and
(iv) nominal phase-dependent speeds.
In this study, the cruise layer is defined as a fixed absolute altitude of \SI{200}{m} above mean sea level (MSL). Because the altitude is defined in MSL, the height above ground may vary with terrain and pad elevation. This is a deliberate benchmark abstraction used to evaluate the computational behavior of the planning stack under controlled conditions. Accordingly, local terrain variation, obstacle constraints, and dynamic weather effects are not modeled in this stage of the evaluation, and the resulting trajectories should be interpreted as benchmark trajectories rather than operational flight plans.

\paragraph{Nominal waypoint synthesis.}
At a high level, the waypoint generator converts the lateral A* route into a
small set of phase-labeled waypoints representing departure, climb, cruise,
descent, and landing.
The lateral corridor path \(P\) is converted into a compact waypoint sequence
\(W=[w_1,\dots,w_m]\) that explicitly represents the take-off, climb, cruise,
descent, and landing phases. The first and last waypoints coincide with the
origin and destination pad locations, respectively. Between these endpoints,
\texttt{acc-trajcalc} inserts a small number of additional waypoints to shape the
vertical profile and to provide CFMS with a phase-structured route definition.

For multi-node routes, intermediate corridor nodes are retained as cruise-route
reference points. Their lateral coordinates follow the graph path, while their
altitudes are assigned according to the nominal climb profile up to the fixed
cruise layer. For direct origin--destination connections, a reduced waypoint
set is generated, including an intermediate point on the connecting geodesic to
support smooth transition between departure and arrival phases.

Each waypoint is annotated with a nominal flight phase and an associated target
speed. The phase labels distinguish terminal segments near the vertiports from
the cruise/transition segments along the corridor network. Nominal speeds are
derived from vehicle performance limits scaled by a safety factor to preserve
flyability margins.

\paragraph{Terminal procedure construction.}
The departure and arrival portions of \(W\) are generated using a simplified
vertiport terminal-procedure model. At departure, the synthesized waypoint
sequence begins at the pad elevation, continues with a vertical climb segment,
and then transitions into the outbound corridor direction. At arrival, the
sequence is constructed analogously in reverse, guiding the vehicle from the
inbound corridor segment toward the destination pad through a constrained
descent geometry.

To keep these terminal segments geometrically structured, the waypoint generator
uses a parameterized obstacle-free-volume, referred to here as a
``verticone,'' around each vertiport. Intuitively, the verticone defines a
protected three-dimensional terminal volume above and around the pad: close to
the pad the motion is mostly vertical, whereas at greater height the allowable
lateral footprint becomes wider. In this study, these geometry
parameters are taken from the EASA Prototype Technical Specifications for the
Design of VFR Vertiports (PTS-VPT-DSN), in particular the vertical take-off and
landing procedure and OFV definitions in Chapter~D, Subpart~2~\cite{easa2022_pts_vpt_dsn}. The synthesized terminal waypoints are chosen
so that departure and arrival remain inside this admissible volume.

Here, \(h_1\) denotes the lower terminal reference height, \(h_2\) the upper
terminal reference height, LDP the landing decision point, OFV the
obstacle-free volume, and OLS the outer limitation surface.
For departure, the vehicle first climbs vertically from the pad to a lower
terminal height \(h_1\). It then transitions toward a second reference height
\(h_2\) while aligning with the outbound corridor direction. If enough
horizontal distance is available, the climb can be extended farther along the
same centerline toward the outer limitation surface at up to \SI{152}{m} above
the pad, limited by the cruise altitude. For arrival, the logic is reversed:
the trajectory may first pass through a top-of-descent point on the inbound
centerline, then enters the OFV at \(h_2\), transitions to an over-pad waypoint
at \(h_1\), passes the LDP reference height, and finishes at the pad.

This procedure produces a sparse but operationally meaningful nominal waypoint
set that preserves the corridor-based lateral route while explicitly modeling
the terminal altitude transitions required before and after cruise flight.
Extended procedure-level variants such as balked landing, holding, or
energy-constrained divert trajectories exist in the implementation, but they
are outside the scope of the present evaluation.

\paragraph{CFMS execution and subsampling.}
The annotated waypoint list is converted into the CFMS input format (WGS-84 coordinates, altitudes, phase and speed fields) and passed to CFMS to produce a dense, time-parameterized trajectory \(\mathcal{T}\) satisfying the kinematic profile~\cite{morscheck2021cfms}.
To reduce payload size for downstream processing and visualization, the output representation is uniformly subsampled to a fixed maximum of 200 trajectory points while retaining the first and last samples. The trajectory generation runtimes reported in the evaluation are measured for CFMS invocation and dense trajectory generation before this subsampling step; the subsampling and visualization-export processing are accounted for separately and are not included in the planner-stage latency decomposition.

% Optional (only if you have a clear pipeline figure from thesis):
% Copy to figures/traj_pipeline.png and uncomment:
% \begin{figure}[!htbp]
%   \centering
%   \includegraphics[width=0.90\linewidth]{figures/traj_pipeline.png}
%   \caption{Trajectory generation pipeline from corridor route to time-parameterized 4D trajectory (CFMS).}
%   \label{fig:traj_pipeline}
% \end{figure}

\subsection{Deconfliction Service Using NDMap}
\label{sec:conflict_detection}

Conflict detection and resolution are performed by the separate \texttt{acc-trajdeconf} service. The service uses NDMap~\cite{NDMap}, configured over the 4D state space \((\lambda,\varphi,H,t)\) (longitude, latitude, altitude, time), to index previously accepted trajectories and identify segments of an inspected candidate trajectory that violate separation minima with respect to the active set. This keeps conflict-management state and NDMap-specific data structures outside the \texttt{acc-trajcalc} path-planning service.

\paragraph{Configuration and output.}
Separation requirements are parameterized by minimum spacing thresholds in each dimension (horizontal, vertical, and temporal).
For each detected conflict, NDMap returns one or more time intervals \([t_{\text{start}}, t_{\text{end}}]\) indicating when the candidate violates separation with any active trajectory. These intervals are used as the primary input to the resolution loop within \texttt{acc-trajdeconf}.

\begin{figure}[!htbp]
\centering
\begin{tikzpicture}[
  node distance=8mm,
  box/.style={
    draw,
    rounded corners,
    align=center,
    inner sep=4pt,
    minimum width=42mm,
    minimum height=8mm
  },
  arrow/.style={-Latex, thick}
]

\node[box] (fetch) {Fetch topology (Webserver)\\+ active traffic (USSP)};
\node[box, below=of fetch] (astar) {\texttt{acc-trajcalc}:\\Global routing (A*)};
\node[box, below=of astar] (traj) {\texttt{acc-trajcalc}:\\4D trajectory generation (CFMS)};
\node[box, below=of traj] (deconf) {\texttt{acc-trajdeconf}:\\Conflict check (NDMap)};
\node[box, below=of deconf] (submit) {Submit conflict-free plan to USSP};

\node[box, right=20mm of deconf] (penalize) {\texttt{acc-trajdeconf}:\\Penalize edges\\and request re-plan};

% Main vertical flow
\draw[arrow] (fetch) -- (astar);
\draw[arrow] (astar) -- (traj);
\draw[arrow] (traj) -- (deconf);
\draw[arrow] (deconf) -- node[right, font=\small] {no conflict} (submit);

% Conflict loop
\draw[arrow] (deconf.east) -- node[above, font=\small] {conflict} (penalize.west);

% Route return path safely below the boxes (no crossing)
\coordinate (retA) at ($(penalize.south)+(0,-6mm)$);
\coordinate (retB) at ($(astar.south)+(0,-6mm)$);

% Return path: go down, go left under all boxes, then go up to A*
\draw[arrow]
  (penalize.north) --
  ([yshift=27mm]penalize.north) --
  (astar.east);
  
\end{tikzpicture}
\caption{Operational planning loop. Candidate trajectories are generated by acc-trajcalc; conflicts are detected and resolved by acc-trajdeconf, which can trigger lateral re-routing and re-execution of A*.}
\label{fig:planning_loop}
\end{figure}

\subsection{Conflict Resolution by Lateral Re-routing}
\label{sec:conflict_resolution}

Conflict resolution is implemented in \texttt{acc-trajdeconf} using \emph{lateral re-routing} on the corridor graph. The procedure is iterative:
(1) request or receive a candidate route from \texttt{acc-trajcalc} using A*,
(2) generate the corresponding 4D trajectory through \texttt{acc-trajcalc} and perform conflict detection in \texttt{acc-trajdeconf} using NDMap,
(3) if conflicts are present, have \texttt{acc-trajdeconf} penalize corridor edges belonging to the conflicting candidate route by assigning them a very large cost (effectively removing them from consideration),
and (4) request re-execution of A* on the modified graph.
The loop terminates when NDMap reports no conflicts or when no feasible alternative route exists.

NDMap, as a trajectory-conflict management concept, can support several resolution dimensions, including changes in time, direction, and altitude~\cite{NDMap}. In the present implementation, \texttt{acc-trajdeconf} uses NDMap for conflict detection and applies resolution at the service level. The current resolution mechanism is restricted to lateral re-routing in the corridor graph; no departure-time adjustment, speed-based temporal shift, or altitude-layer change is applied.

This strategy is intentionally conservative and simple: it preserves the overall system modularity and enables scalable online admission control while illustrating how 4D deconfliction can be integrated as a dedicated microservice. More expressive resolution strategies based on temporal or vertical deconfliction are left for future work.

% Optional (use a REAL before/after figure if you have it):
% Copy to figures/conflict_before_after.png and uncomment:
% \begin{figure}[!htbp]
%   \centering
%   \includegraphics[width=0.95\linewidth]{figures/conflict_before_after.png}
%   \caption{Example of conflict handling via lateral re-routing: initial candidate trajectory conflicts with active traffic; the planner penalizes conflicting edges and produces an alternative conflict-free route.}
%   \label{fig:conflict_before_after}
% \end{figure}
